\documentclass{article}

%-------------------------%
% PREAMBUŁA --------------%
%-------------------------%
\usepackage[utf8]{inputenc}
\usepackage[OT4]{polski}
\usepackage{caption}
\usepackage{amsmath}
\usepackage{tikz}
\usepackage{subcaption}
\usepackage{tabularx}
\usepackage{caption}
\usepackage{array}
\usepackage{hyperref}
\usepackage{graphicx}
\usepackage{fancyhdr}
\usepackage[justification=centering]{caption}
\usepackage[margin=60pt]{geometry}
\usepackage{longtable}



%\usepackage[T1]{fontenc}
%\usepackage{amsthm}
%\theoremstyle{plain}
%\usepackage{enumitem}
%\captionsetup[table]{skip=-4pt}
%\usepackage{tabulary}
%\usepackage{sidecap}
%\usepackage{gensymb}
%\usepackage{wrapfig}
%\usepackage{lastpage}
%\usepackage{gensymb}
%\usepackage{lipsum}
%\pagestyle{fancy}


\newlength{\RoundedBoxWidth}
\newsavebox{\GrayRoundedBox}
\newenvironment{GrayBox}[1][\dimexpr\textwidth-4.5ex]%
   {\setlength{\RoundedBoxWidth}{\dimexpr#1}
    \begin{lrbox}{\GrayRoundedBox}
       \begin{minipage}{\RoundedBoxWidth}}%
   {   \end{minipage}
    \end{lrbox}
    \begin{center}
    \begin{tikzpicture}%
       \draw node[draw=black,fill=black!1,rounded corners,%
             inner sep=2ex,text width=\RoundedBoxWidth]%
             {\usebox{\GrayRoundedBox}};
    \end{tikzpicture}
    \end{center}}

%------- Typy kolumn do ustawiania szerokości ------- %
\newcolumntype{L}[1]{>{\raggedright\let\newline\\\arraybackslash\hspace{0pt}}m{#1}}
\newcolumntype{C}[1]{>{\centering\let\newline\\\arraybackslash\hspace{0pt}}m{#1}}
\newcolumntype{R}[1]{>{\raggedleft\let\newline\\\arraybackslash\hspace{0pt}}m{#1}}








\begin{document}

%-------------------------%
%Tabela nagłówkowa -------%
%-------------------------%

\begin{table}[h]
\begin{tabular}{|l|l|l|l|l|l|}
\hline
\begin{tabular}[c]{@{}l@{}}
Wydział:\\ WFiIS\end{tabular} &
\multicolumn{2}{l|}{\begin{tabular}[c]{@{}l@{}}
Imię i nazwisko:\\ 1. Axel Zuziak\\ 2. Marcin Węglarz \end{tabular}} 
& Rok \textbf{II}         
& Grupa \textbf{B}          
& Zespół \textbf{03}      \\ \hline
\textbf{\begin{tabular}[c]{@{}l@{}}
LABOLATORIUM\\ TECHNIK \\ JĄDROWYCH\end{tabular}} &
\multicolumn{4}{l|}{ Temat:\textbf{ Statystyczny charakter rozpadów promieniotwórczych.}
	 }                                                                                                                       & \multicolumn{1}{c|}{\begin{tabular}[c]{@{}c@{}}
Nr ćwiczenia\\ \textbf{1+9}
\end{tabular}} \\ \hline
\begin{tabular}[c]{@{}c@{}}
Data wykonania:\\ 04.03.2015
\end{tabular} &
\begin{tabular}[c]{@{}c@{}}
 Data oddania:\\ 18.03.2015
\end{tabular} &
\begin{tabular}[c]{@{}c@{}}
Zwrot do poprawy: \\ 
\end{tabular}                                 
& 
\begin{tabular}[c]{@{}c@{}}
Data oddania: \\ 
\end{tabular}&
\begin{tabular}[c]{@{}c@{}}
Data zaliczenia:\\
\end{tabular}& 
\begin{tabular}[c]{@{}c@{}}
OCENA:\\
\end{tabular}       
\\ & & & & & \\ \hline
\end{tabular}
\end{table}
%-----Koniec tabeli------

%%%%%%%%%%%% Wstęp %%%%%%%%%%%%
\section{Wstęp teoretyczny}
\section{Aparatura i wykonanie ćwiczenia}
\begin{itemize}
	\item \textbf{Neutronowy miernik wodoru}
	\item \textbf{Wzmacniacz impulsowy}
	\item \textbf{Zasilacz wysokiego napięcia}
	\item \textbf{Analizator amplitudy}
	\item \textbf{Przelicznik}
\end{itemize}

Ćwiczenie rozpoczęto od zapoznania się z aparaturą pomiarową. Przed przystąpieniem do wykonywania pomiarów zmierzono wszystkie wymiary geometryczne zarówno próbek wzorcowych: $W1,W2,...W7$ jak i tych, dla których wyznaczano gęstość wodoru: $P1,P2,...P7$. Zanotowano wagi poszczególnych próbek. Wykonano pomiary przy pomocy neutronowego miernika wodoru dla wszystkich próbek.

\section{Wyniki pomiarów i obliczenia.}

\subsection{Wyznaczenie objętości i gęstości badanych próbek.}
Przy pomocy linijki zmierzono wymiary geometryczne badanych próbek wyniki pomiarów i obliczeń przedstawiono w tabeli \ref{geometryczne}.

%%%%%%%%%%%% BIBLIOGRAFIA %%%%%%%%%%%%

\begin{thebibliography}{9}
	
	
	\bibitem{1}
	\url{http://nucleardata.nuclear.lu.se/toi/}

	
	\bibitem{DzK}
	B. Dziunikowski, S.J. Kalita
	\emph{Ćwiczenia laboratoryjne z jądrowych metod pomiarowcyh}, Wydawnictwa AGH, Kraków 1995
	
	\bibitem{przybycien}
	Mariusz Przybycień \emph{Tablice Statystyczne}, dostęp on-line\\
	\url{http://home.agh.edu.pl/~mariuszp/wfiis_stat/tablice_ps_wir.pdf}
\end{thebibliography}
\vspace{2cm}








\end{document}





%\textsc{http://nucleardata.nuclear.lu.se/toi/}
%do bibliografii






